\documentclass[11pt,a4paper]{article}
\usepackage[utf8]{inputenc}
\usepackage[T1]{fontenc}
\usepackage[catalan]{babel}
\usepackage{lmodern}
\usepackage[sfdefault,lining]{FiraSans}
\usepackage[a4paper,top=2.3cm,bottom=2.2cm,left=1.8cm,right=1.8cm,headheight=1.3cm,headsep=0.4cm,footskip=1cm]{geometry}
\usepackage{amsmath,amssymb}
\allowdisplaybreaks
\usepackage{icomma}
\usepackage{graphicx}
\usepackage[table]{xcolor}
\usepackage{tikz}
\usepackage{fontawesome5}
\usepackage{ccicons}
\usepackage{fancyhdr}
\usepackage{lastpage}
\usepackage{tasks}
\usepackage{needspace}
\usepackage[most]{tcolorbox}
\usepackage{array}
\usepackage[hidelinks,pdftitle={Solucions de la fitxa de la unitat 5: Límits i continuïtat. 2n de Batxillerat MII},pdfsubject={Matemàtiques II},pdfkeywords={límits, indeterminacions, infinitèsims equivalents, límits trigonomètrics, paràmetres, continuïtat, funcions a trossos, teorema de Bolzano, asímptotes, PAU}]{hyperref}

\definecolor{navy}{HTML}{1B2748}
\definecolor{accent}{HTML}{E8603A}
\definecolor{teal}{HTML}{1F9E90}
\definecolor{soft}{HTML}{EEF1F8}
\definecolor{muted}{HTML}{5D6685}

\setlength{\parindent}{0pt}
\setlength{\parskip}{3pt}
\settasks{label=\alph*), label-format=\bfseries\color{navy}, label-width=1.3em, item-indent=1.8em, column-sep=0.6em, after-item-skip=5pt, before-skip=3pt, after-skip=4pt}

% ---------- capçalera i peu ----------
\newcommand{\llicencia}{\href{https://creativecommons.org/licenses/by-nc-sa/4.0/deed.ca}{\ccbyncsa}\ \ Llicència Creative Commons Reconeixement-NoComercial-CompartirIgual 4.0 Internacional (CC BY-NC-SA 4.0)}
\pagestyle{fancy}
\fancyhf{}
\fancyhead[L]{\raisebox{-0.35\height}{\includegraphics[height=0.95cm]{logo-gva.png}}}
\fancyhead[C]{\small\color{muted}Solucions de la fitxa · Unitat 5: Límits i continuïtat · 2n Batxillerat MII}
\fancyhead[R]{\raisebox{-0.35\height}{\includegraphics[height=0.95cm]{logo-vermellar.png}}}
\fancyfoot[L]{\scriptsize\color{muted}\llicencia}
\fancyfoot[R]{\small\color{muted}\thepage}
\renewcommand{\headrulewidth}{0.4pt}
\renewcommand{\headrule}{\color{muted!40}\hrule height \headrulewidth}
\fancypagestyle{primera}{\fancyhf{}\renewcommand{\headrulewidth}{0pt}
  \fancyfoot[L]{\scriptsize\color{muted}\llicencia}
  \fancyfoot[R]{\small\color{muted}\thepage}}

% ---------- nivells i exercicis ----------
\newcommand{\estrelles}[1]{{\color{accent}\ifcase#1\or\faStar\faStar[regular]\faStar[regular]\or\faStar\faStar\faStar[regular]\or\faStar\faStar\faStar\fi}}
\newcounter{exer}
\newcommand{\ex}[1]{\par\needspace{4\baselineskip}\medskip\refstepcounter{exer}\noindent{\color{navy}\bfseries\theexer.}\ \ }

\newcommand{\seccio}[2]{\par\needspace{10\baselineskip}\bigskip
  \begin{tcolorbox}[colback=navy,colframe=navy,coltext=white,arc=2mm,boxrule=0pt,left=3mm,right=3mm,top=1.2mm,bottom=1.2mm]
  {\large\bfseries #1\quad #2}
  \end{tcolorbox}\vspace{-2pt}}


\setlength{\emergencystretch}{3em}

% quadre de recordatori (només a la fitxa d'exercicis)
\newcommand{\recorda}[1]{\par\smallskip\begin{tcolorbox}[colback=accent!7,colframe=accent!55,arc=2mm,boxrule=0.5pt,left=2.5mm,right=2.5mm,top=1mm,bottom=1mm]\small{\color{accent}\faLightbulb[regular]}\ \textbf{Recorda:}\ #1\end{tcolorbox}}
% etiqueta dels exercicis de la PAU d'altres comunitats
\newcommand{\pau}[1]{{\color{teal}\small\bfseries\faUniversity\ PAU · #1.}\ }
\newcommand{\res}[1]{\boxed{#1}}\newcommand{\rest}[1]{\fbox{#1}}

% =====================================================================
% Capçalera FentMates de les fitxes: un color per a cada curs.
% Els colors són els de la pàgina inicial (LEVELS a assets/js/materials.js);
% el logo ve de _eines/logo/pdf/ (python3 _eines/logo/genera.py).
%
% Ús, al final del preàmbul de la unitat:
%   % =====================================================================
% Capçalera FentMates de les fitxes: un color per a cada curs.
% Els colors són els de la pàgina inicial (LEVELS a assets/js/materials.js);
% el logo ve de _eines/logo/pdf/ (python3 _eines/logo/genera.py).
%
% Ús, al final del preàmbul de la unitat:
%   % =====================================================================
% Capçalera FentMates de les fitxes: un color per a cada curs.
% Els colors són els de la pàgina inicial (LEVELS a assets/js/materials.js);
% el logo ve de _eines/logo/pdf/ (python3 _eines/logo/genera.py).
%
% Ús, al final del preàmbul de la unitat:
%   \input{../comu/fentmates.tex}
%   \fmcurs{1eso}
% i, al començament del document:
%   \capcalera{Matemàtiques}{1r ESO}{Fitxa d'exercicis · Unitat 2}{Divisibilitat}
% Opcional abans de \input: \def\fmneed{14} (espai mínim abans de cada secció).
% =====================================================================
\makeatletter
\def\fm@logodir{../../logo/pdf}
\providecommand{\fmneed}{10}
\def\fm@cursnom{}

% clau del curs, color i color del text damunt del color
\newcommand{\fm@curs}[3]{\@namedef{fm@color@#1}{#2}\@namedef{fm@text@#1}{#3}}
\fm@curs{1eso}{E6B84A}{navy}
\fm@curs{2eso}{3B82F6}{white}
\fm@curs{3eso}{FF7A50}{white}
\fm@curs{4eso-a}{A78BFA}{white}
\fm@curs{4eso-b}{6D28D9}{white}
\fm@curs{1batx-mi}{9B1C1C}{white}
\fm@curs{1batx-ccss}{EF6B6B}{white}
\fm@curs{2batx-mii}{1E6B47}{white}
\fm@curs{2batx-ccss}{2BB3A3}{white}

\newcommand{\fmcurs}[1]{%
  \@ifundefined{fm@color@#1}{\PackageError{fentmates}{Curs desconegut: #1}{}}{}%
  \edef\fm@tmp{\noexpand\definecolor{curs}{HTML}{\@nameuse{fm@color@#1}}%
               \noexpand\colorlet{curstext}{\@nameuse{fm@text@#1}}}\fm@tmp
  \colorlet{cursfosc}{curs!72!black}%
  \def\fm@logo{\fm@logodir/vertical-#1.pdf}%
  \def\fm@marca{\fm@logodir/marca-#1.pdf}}

% ---------- capçalera de la primera pàgina ----------
% #1 matèria (petit), #2 curs, #3 tipus de document i unitat, #4 títol
\newcommand{\capcalera}[4]{%
  \gdef\fm@cursnom{#2}%
  \par\noindent
  \begin{tcolorbox}[colback=curs,colframe=curs,boxrule=0pt,arc=2.2mm,width=3.5cm,height=2.15cm,
      valign=center,left=3mm,right=2mm,top=1mm,bottom=1mm,nobeforeafter,box align=center]
    \color{curstext}{\footnotesize #1}\par\vspace{1pt}{\fontsize{19}{21}\selectfont\bfseries #2\par}%
  \end{tcolorbox}\hspace{5mm}%
  \begin{minipage}[c]{\dimexpr\linewidth-3.5cm-5mm-2.9cm\relax}
    {\color{cursfosc}\small\bfseries\MakeUppercase{#3}}\par\vspace{2pt}
    {\color{navy}\fontsize{20}{23}\selectfont\bfseries #4\par}
  \end{minipage}\hfill
  \begin{minipage}[c]{2.6cm}\raggedleft\includegraphics[width=2.6cm]{\fm@logo}\end{minipage}\par
  \vspace{2pt}}

% ---------- barres de secció del color del curs ----------
\renewcommand{\seccio}[2]{\par\needspace{\fmneed\baselineskip}\bigskip
  \begin{tcolorbox}[colback=curs,colframe=curs,coltext=curstext,arc=2mm,boxrule=0pt,left=3mm,right=3mm,top=1.2mm,bottom=1.2mm]
  {\large\bfseries #1\quad #2}
  \end{tcolorbox}\vspace{-2pt}}

% ---------- estrelles de dificultat del color del curs (a Batxillerat no se'n posen) ----------
\@ifundefined{estrelles}{}{%
  \renewcommand{\estrelles}[1]{{\color{curs}\ifcase#1\or\faStar\faStar[regular]\faStar[regular]\or\faStar\faStar\faStar[regular]\or\faStar\faStar\faStar\fi}}}

% ---------- capçalera i peu de les pàgines ----------
% primera pàgina: logos institucionals al peu
\fancypagestyle{primera}{\fancyhf{}\renewcommand{\headrulewidth}{0pt}
  \fancyfoot[L]{\raisebox{-0.32\height}{\includegraphics[height=0.78cm]{logo-gva.png}}\hspace{5mm}%
    \raisebox{-0.32\height}{\includegraphics[height=0.78cm]{logo-vermellar.png}}}
  \fancyfoot[R]{\scriptsize\color{muted}\href{https://creativecommons.org/licenses/by-nc-sa/4.0/deed.ca}{\ccbyncsa}\ \ CC BY-NC-SA 4.0 · fentmates.org\qquad\small\thepage}}
% la resta: el símbol a l'esquerra i el curs a la dreta, amb el títol al centre com abans
\fancyhead[L]{\raisebox{-0.3\height}{\includegraphics[height=0.62cm]{\fm@marca}}}
\fancyhead[R]{\small\bfseries\color{cursfosc}\fm@cursnom}
\renewcommand{\headrule}{\color{curs!45}\hrule height \headrulewidth}
\makeatother

%   \fmcurs{1eso}
% i, al començament del document:
%   \capcalera{Matemàtiques}{1r ESO}{Fitxa d'exercicis · Unitat 2}{Divisibilitat}
% Opcional abans de \input: \def\fmneed{14} (espai mínim abans de cada secció).
% =====================================================================
\makeatletter
\def\fm@logodir{../../logo/pdf}
\providecommand{\fmneed}{10}
\def\fm@cursnom{}

% clau del curs, color i color del text damunt del color
\newcommand{\fm@curs}[3]{\@namedef{fm@color@#1}{#2}\@namedef{fm@text@#1}{#3}}
\fm@curs{1eso}{E6B84A}{navy}
\fm@curs{2eso}{3B82F6}{white}
\fm@curs{3eso}{FF7A50}{white}
\fm@curs{4eso-a}{A78BFA}{white}
\fm@curs{4eso-b}{6D28D9}{white}
\fm@curs{1batx-mi}{9B1C1C}{white}
\fm@curs{1batx-ccss}{EF6B6B}{white}
\fm@curs{2batx-mii}{1E6B47}{white}
\fm@curs{2batx-ccss}{2BB3A3}{white}

\newcommand{\fmcurs}[1]{%
  \@ifundefined{fm@color@#1}{\PackageError{fentmates}{Curs desconegut: #1}{}}{}%
  \edef\fm@tmp{\noexpand\definecolor{curs}{HTML}{\@nameuse{fm@color@#1}}%
               \noexpand\colorlet{curstext}{\@nameuse{fm@text@#1}}}\fm@tmp
  \colorlet{cursfosc}{curs!72!black}%
  \def\fm@logo{\fm@logodir/vertical-#1.pdf}%
  \def\fm@marca{\fm@logodir/marca-#1.pdf}}

% ---------- capçalera de la primera pàgina ----------
% #1 matèria (petit), #2 curs, #3 tipus de document i unitat, #4 títol
\newcommand{\capcalera}[4]{%
  \gdef\fm@cursnom{#2}%
  \par\noindent
  \begin{tcolorbox}[colback=curs,colframe=curs,boxrule=0pt,arc=2.2mm,width=3.5cm,height=2.15cm,
      valign=center,left=3mm,right=2mm,top=1mm,bottom=1mm,nobeforeafter,box align=center]
    \color{curstext}{\footnotesize #1}\par\vspace{1pt}{\fontsize{19}{21}\selectfont\bfseries #2\par}%
  \end{tcolorbox}\hspace{5mm}%
  \begin{minipage}[c]{\dimexpr\linewidth-3.5cm-5mm-2.9cm\relax}
    {\color{cursfosc}\small\bfseries\MakeUppercase{#3}}\par\vspace{2pt}
    {\color{navy}\fontsize{20}{23}\selectfont\bfseries #4\par}
  \end{minipage}\hfill
  \begin{minipage}[c]{2.6cm}\raggedleft\includegraphics[width=2.6cm]{\fm@logo}\end{minipage}\par
  \vspace{2pt}}

% ---------- barres de secció del color del curs ----------
\renewcommand{\seccio}[2]{\par\needspace{\fmneed\baselineskip}\bigskip
  \begin{tcolorbox}[colback=curs,colframe=curs,coltext=curstext,arc=2mm,boxrule=0pt,left=3mm,right=3mm,top=1.2mm,bottom=1.2mm]
  {\large\bfseries #1\quad #2}
  \end{tcolorbox}\vspace{-2pt}}

% ---------- estrelles de dificultat del color del curs (a Batxillerat no se'n posen) ----------
\@ifundefined{estrelles}{}{%
  \renewcommand{\estrelles}[1]{{\color{curs}\ifcase#1\or\faStar\faStar[regular]\faStar[regular]\or\faStar\faStar\faStar[regular]\or\faStar\faStar\faStar\fi}}}

% ---------- capçalera i peu de les pàgines ----------
% primera pàgina: logos institucionals al peu
\fancypagestyle{primera}{\fancyhf{}\renewcommand{\headrulewidth}{0pt}
  \fancyfoot[L]{\raisebox{-0.32\height}{\includegraphics[height=0.78cm]{logo-gva.png}}\hspace{5mm}%
    \raisebox{-0.32\height}{\includegraphics[height=0.78cm]{logo-vermellar.png}}}
  \fancyfoot[R]{\scriptsize\color{muted}\href{https://creativecommons.org/licenses/by-nc-sa/4.0/deed.ca}{\ccbyncsa}\ \ CC BY-NC-SA 4.0 · fentmates.org\qquad\small\thepage}}
% la resta: el símbol a l'esquerra i el curs a la dreta, amb el títol al centre com abans
\fancyhead[L]{\raisebox{-0.3\height}{\includegraphics[height=0.62cm]{\fm@marca}}}
\fancyhead[R]{\small\bfseries\color{cursfosc}\fm@cursnom}
\renewcommand{\headrule}{\color{curs!45}\hrule height \headrulewidth}
\makeatother

%   \fmcurs{1eso}
% i, al començament del document:
%   \capcalera{Matemàtiques}{1r ESO}{Fitxa d'exercicis · Unitat 2}{Divisibilitat}
% Opcional abans de \input: \def\fmneed{14} (espai mínim abans de cada secció).
% =====================================================================
\makeatletter
\def\fm@logodir{../../logo/pdf}
\providecommand{\fmneed}{10}
\def\fm@cursnom{}

% clau del curs, color i color del text damunt del color
\newcommand{\fm@curs}[3]{\@namedef{fm@color@#1}{#2}\@namedef{fm@text@#1}{#3}}
\fm@curs{1eso}{E6B84A}{navy}
\fm@curs{2eso}{3B82F6}{white}
\fm@curs{3eso}{FF7A50}{white}
\fm@curs{4eso-a}{A78BFA}{white}
\fm@curs{4eso-b}{6D28D9}{white}
\fm@curs{1batx-mi}{9B1C1C}{white}
\fm@curs{1batx-ccss}{EF6B6B}{white}
\fm@curs{2batx-mii}{1E6B47}{white}
\fm@curs{2batx-ccss}{2BB3A3}{white}

\newcommand{\fmcurs}[1]{%
  \@ifundefined{fm@color@#1}{\PackageError{fentmates}{Curs desconegut: #1}{}}{}%
  \edef\fm@tmp{\noexpand\definecolor{curs}{HTML}{\@nameuse{fm@color@#1}}%
               \noexpand\colorlet{curstext}{\@nameuse{fm@text@#1}}}\fm@tmp
  \colorlet{cursfosc}{curs!72!black}%
  \def\fm@logo{\fm@logodir/vertical-#1.pdf}%
  \def\fm@marca{\fm@logodir/marca-#1.pdf}}

% ---------- capçalera de la primera pàgina ----------
% #1 matèria (petit), #2 curs, #3 tipus de document i unitat, #4 títol
\newcommand{\capcalera}[4]{%
  \gdef\fm@cursnom{#2}%
  \par\noindent
  \begin{tcolorbox}[colback=curs,colframe=curs,boxrule=0pt,arc=2.2mm,width=3.5cm,height=2.15cm,
      valign=center,left=3mm,right=2mm,top=1mm,bottom=1mm,nobeforeafter,box align=center]
    \color{curstext}{\footnotesize #1}\par\vspace{1pt}{\fontsize{19}{21}\selectfont\bfseries #2\par}%
  \end{tcolorbox}\hspace{5mm}%
  \begin{minipage}[c]{\dimexpr\linewidth-3.5cm-5mm-2.9cm\relax}
    {\color{cursfosc}\small\bfseries\MakeUppercase{#3}}\par\vspace{2pt}
    {\color{navy}\fontsize{20}{23}\selectfont\bfseries #4\par}
  \end{minipage}\hfill
  \begin{minipage}[c]{2.6cm}\raggedleft\includegraphics[width=2.6cm]{\fm@logo}\end{minipage}\par
  \vspace{2pt}}

% ---------- barres de secció del color del curs ----------
\renewcommand{\seccio}[2]{\par\needspace{\fmneed\baselineskip}\bigskip
  \begin{tcolorbox}[colback=curs,colframe=curs,coltext=curstext,arc=2mm,boxrule=0pt,left=3mm,right=3mm,top=1.2mm,bottom=1.2mm]
  {\large\bfseries #1\quad #2}
  \end{tcolorbox}\vspace{-2pt}}

% ---------- estrelles de dificultat del color del curs (a Batxillerat no se'n posen) ----------
\@ifundefined{estrelles}{}{%
  \renewcommand{\estrelles}[1]{{\color{curs}\ifcase#1\or\faStar\faStar[regular]\faStar[regular]\or\faStar\faStar\faStar[regular]\or\faStar\faStar\faStar\fi}}}

% ---------- capçalera i peu de les pàgines ----------
% primera pàgina: logos institucionals al peu
\fancypagestyle{primera}{\fancyhf{}\renewcommand{\headrulewidth}{0pt}
  \fancyfoot[L]{\raisebox{-0.32\height}{\includegraphics[height=0.78cm]{logo-gva.png}}\hspace{5mm}%
    \raisebox{-0.32\height}{\includegraphics[height=0.78cm]{logo-vermellar.png}}}
  \fancyfoot[R]{\scriptsize\color{muted}\href{https://creativecommons.org/licenses/by-nc-sa/4.0/deed.ca}{\ccbyncsa}\ \ CC BY-NC-SA 4.0 · fentmates.org\qquad\small\thepage}}
% la resta: el símbol a l'esquerra i el curs a la dreta, amb el títol al centre com abans
\fancyhead[L]{\raisebox{-0.3\height}{\includegraphics[height=0.62cm]{\fm@marca}}}
\fancyhead[R]{\small\bfseries\color{cursfosc}\fm@cursnom}
\renewcommand{\headrule}{\color{curs!45}\hrule height \headrulewidth}
\makeatother

\fmcurs{2batx-mii}

\begin{document}
\thispagestyle{primera}

\capcalera{Matemàtiques II}{2n Batx}{Solucions de la fitxa · Unitat 5}{\raggedright Límits i continuïtat}
\vspace{8pt}

\begin{tcolorbox}[colback=soft,colframe=soft,arc=2mm,left=3mm,right=3mm,top=1.5mm,bottom=1.5mm]
\small Intenta resoldre cada exercici abans de mirar-ne la solució. Si no t'ix el mateix resultat, repassa els passos o comprova'l amb la calculadora de límits dels apunts interactius de la unitat. Cap d'aquests límits necessita derivades: n'hi ha prou amb operar, el conjugat, la comparació d'infinits, la fórmula de $1^\infty$ i els infinitèsims equivalents.
\end{tcolorbox}

% =====================================================
\seccio{1}{Límits a l'infinit}

\ex{2} \begin{tasks}(3)
  \task $\lim 2x^5=\res{-\infty}$
  \task $\frac42=\res{2}$
  \task $\lim(-3x)=\res{+\infty}$
  \task $\res{0}$
  \task $\lim\frac{2+5\cdot2^{-x}}{1-3\cdot2^{-x}}=\res{2}$
  \task $\frac{0-0}{0+4}=\res{0}$
\end{tasks}

\ex{2} \begin{tasks}(3)
  \task $\frac{3x}{2x}\to\res{\frac32}$
  \task $\frac{-3x}{2x}\to\res{-\frac32}$
  \task $\frac{3x-x}{2x}\to\res{1}$
  \task $\frac{2x}{x}\to\res{2}$
  \task $\frac{x^{3/2}}{x}=\sqrt x\to\res{+\infty}$
  \task $\frac{-x}{3x}\to\res{-\frac13}$
\end{tasks}

\ex{2} $\res{\ln x\ll(\ln x)^3\ll\sqrt x\ll x^3\ll x^{10}\ll2^x\ll e^x}$ (qualsevol potència de $\ln x$ és menor que qualsevol potència de $x$).\begin{tasks}(3)
  \task $\res{0}$
  \task $\res{0}$
  \task $\frac{2^x}{3\cdot 2^x}\to\res{\frac13}$
  \task $\res{-\infty}$
  \task $\frac{3^x}{3^x}\to\res{1}$
  \task $\res{0}$
\end{tasks}

\ex{2} \begin{tasks}(2)
  \task $\lim\ln\frac{2x+1}{x+3}=\res{\ln2}$
  \task $\lim\ln\frac{x^2+1}{x}=\ln(+\infty)=\res{+\infty}$
\end{tasks}

\ex{2} \pau{Andalusia, ordinària 2023}$\displaystyle\lim_{x\to+\infty}\frac{x^2}{e^x+e^{-x}}=\res{0}$: l'exponencial és un infinit d'ordre superior a $x^2$.


% =====================================================
\seccio{2}{Indeterminacions a l'infinit}

\ex{2} \begin{tasks}(2)
  \task $\dfrac{-x}{x+1}\to\res{-1}$
  \task $\dfrac{x^3-x^2}{x^3+x^2+x+1}\to\res{1}$
  \task $\dfrac{8x}{\sqrt{x^2+8x}+x}\to\dfrac{8x}{2x}=\res{4}$
  \task $\dfrac{4x}{\sqrt{x^2+3x}+\sqrt{x^2-x}}\to\res{2}$
  \task $x=-t$: \ $\dfrac{-2t}{\sqrt{t^2-2t}+t}\to\res{-1}$
  \task $x-3x=-2x\to\res{-\infty}$
\end{tasks}

\ex{3} Amb $A=\sqrt[3]{x^3+3x^2}$, $B=x$: \ $\dfrac{A^3-B^3}{A^2+AB+B^2}=\dfrac{3x^2}{A^2+Ax+x^2}\to\dfrac{3x^2}{3x^2}=\res{1}$

\ex{2} \begin{tasks}(2)
  \task $\dfrac{4x}{\sqrt{x^2+4}+x}\to\res{2}$
  \task $\dfrac{2x+1}{\sqrt{x+1}+\sqrt x}\approx\sqrt x\to\res{+\infty}$
  \task $\ln\left(1+\frac2x\right)^x\to\ln e^2=\res{2}$
  \task $\dfrac{x^3}{e^x}\to\res{0}$
\end{tasks}

\ex{2} \begin{tasks}(2)
  \task $e^{\lim\frac{3}{x-1}\cdot2x}=\res{e^6}$
  \task $e^{\lim\frac{-4(x+1)}{2x+3}}=\res{e^{-2}}$
  \task $e^{\lim\frac{4x^2}{x^2-1}}=\res{e^4}$
  \task $3^{+\infty}=\res{+\infty}$ (no és $1^\infty$)
  \task $e^{\lim\frac{-4}{x}\cdot3x}=\res{e^{-12}}$
  \task $e^{\lim\frac{(-x-1)(x^2+1)}{(x^2+2)(x-3)}}=\res{e^{-1}}$
\end{tasks}

\ex{3} \begin{tasks}(2)
  \task $\ln L=\lim\frac{\ln x}{x}=0\Rightarrow L=\res{1}$
  \task $\ln L=\lim\frac{2\ln3+2\ln x}{\ln x}=2\Rightarrow L=\res{e^2}$
\end{tasks}

\ex{2} \pau{Aragó, ordinària 2024}$\dfrac{x^2+5-(x+2)^2}{\sqrt{x^2+5}+x+2}=\dfrac{-4x+1}{\sqrt{x^2+5}+x+2}\to\dfrac{-4x}{2x}=\res{-2}$

\ex{2} \pau{Aragó, ordinària 2022}$\dfrac{3x^2-2-(\sqrt3x+5)^2}{\sqrt{3x^2-2}+\sqrt3x+5}=\dfrac{-10\sqrt3\,x-27}{\sqrt{3x^2-2}+\sqrt3x+5}\to\dfrac{-10\sqrt3\,x}{2\sqrt3\,x}=\res{-5}$

\ex{2} \pau{Navarra, extraordinària 2021}Pel conjugat: $\dfrac{\sqrt x\left(\sqrt{3x^3+2x^2}+\sqrt{3x^3}\right)}{2x^2}\approx\dfrac{\sqrt x\cdot2\sqrt3\,x^{3/2}}{2x^2}\to\res{\sqrt3}$

\ex{2} \pau{Múrcia, ordinària 2024}\begin{tasks}(2)
  \task $\dfrac{18}{\sqrt{x+9}+\sqrt{x-9}}\to\res{0}$
  \task $\res{0}$ (comparació d’infinits)
\end{tasks}

\ex{2} \pau{Castella-la Manxa, ordinària 2023}$1^\infty$: \ $e^{\lim\frac{1}{5x}\cdot x^2}=e^{+\infty}=\res{+\infty}$

\ex{2} \pau{Navarra, extraordinària 2022}$\dfrac{x^2}{x^2+x-1}-1=\dfrac{1-x}{x^2+x-1}$; \ $\lim\dfrac{(1-x)(2x-1)}{x^2+x-1}=-2$: \ $\res{e^{-2}}$

\ex{2} \pau{La Rioja, extraordinària 2023}$\infty^0$: \ $\ln L=\lim\dfrac{\ln(1+x)}{x}=0$ (comparació d'infinits) $\Rightarrow L=\res{1}$

\ex{2} \pau{Illes Balears, extraordinària 2021}\begin{tasks}(1)
  \task $P(0)=\res{15}$ milions
  \task $\displaystyle\lim_{t\to+\infty}P(t)=\res{1}$ milió d'individus
\end{tasks}


% =====================================================
\seccio{3}{Límits en un nombre i indeterminacions}

\ex{2} \begin{tasks}(1)
  \task $\frac{3}{0^+}$ pels dos costats: \ $\res{+\infty}$
  \task $x^3-5x^2+3x+9=(x-3)^2(x+1)$: \ $\dfrac{1}{(x-3)(x+1)}$; \ $\lim_{x\to3^-}=\res{-\infty}$, \ $\lim_{x\to3^+}=\res{+\infty}$: \ $\res{\nexists}$ el límit
\end{tasks}

\ex{2} \begin{tasks}(1)
  \task $x^2+2x+4\to\res{12}$
  \task $\dfrac{(x-1)(x+1)(x^2+1)}{x(x-1)}\to\res{4}$
  \task $\dfrac{(x-2)^3(x+2)}{(x-2)^3}=x+2\to\res{4}$
\end{tasks}

\ex{2} \begin{tasks}(1)
  \task $\dfrac{x-4}{(x-4)(\sqrt{x+5}+3)}\to\res{\frac16}$
  \task $\dfrac{-(x-3)}{(x-3)(x+3)\left(\sqrt{x+1}+\sqrt{2x-2}\right)}\to\dfrac{-1}{6\cdot4}=\res{-\frac{1}{24}}$
  \task $\dfrac{x\left(\sqrt{4+x}+2\right)}{x\left(\sqrt{1+x}+1\right)}\to\res{2}$
\end{tasks}

\ex{2} \begin{tasks}(2)
  \task $\dfrac{x-1}{(x-1)(x+1)}\to\res{\frac12}$
  \task $\dfrac{(x-2)(x+4)}{(x-2)(x+2)}\to\res{\frac32}$
  \task $e^{\lim3x\cdot\frac2x}=\res{e^6}$
  \task $e^{\lim\frac{-1}{2x}}=\res{e^{-1/4}}$
  \task $e^{\lim\frac{3(x-1)}{x-1}}=\res{e^3}$
  \task $e^{\lim\frac{2(x-1)}{x-1}}=\res{e^2}$
\end{tasks}

\ex{2} \pau{Castella-la Manxa, extraordinària 2024}$\dfrac{(x+2)(x^2+1)}{x(x+2)}=\dfrac{x^2+1}{x}\to\res{-\dfrac52}$

\ex{2} \pau{Castella-la Manxa, extraordinària 2023}$\dfrac{(x-3)(x^2+3)}{3(x-3)}\to\dfrac{12}{3}=\res{4}$

\ex{2} \pau{Castella i Lleó, ordinària 2023}Pel conjugat: $\dfrac{x-2}{(x-2)\left(\sqrt{x^3+x-1}+\sqrt{x^3+1}\right)}\to\dfrac{1}{3+3}=\res{\dfrac16}$

\ex{2} \pau{Múrcia, ordinària 2023}$\dfrac{2x}{3x\left(\sqrt{9+x}+\sqrt{9-x}\right)}\to\dfrac{2}{3\cdot6}=\res{\dfrac19}$


% =====================================================
\seccio{4}{Límits trigonomètrics i infinitèsims equivalents}

\ex{2} \begin{tasks}(3)
  \task $\frac{4x}{3x}=\res{\frac43}$
  \task $\frac{x\cdot2x}{x^2/2}=\res{4}$
  \task $\frac{x^2}{9x^2}=\res{\frac19}$
  \task $\frac{3x}{5x}=\res{\frac35}$
  \task $\frac{(2x)^2/2}{x^2}=\res{2}$
  \task $\frac{2x}{3x}=\res{\frac23}$
\end{tasks}

\ex{2} \begin{tasks}(2)
  \task $\frac{e^{-x}(e^{2x}-1)}{\sin x}\sim\frac{2x}{x}\to\res{2}$
  \task $t=x-\frac\pi2$: $\frac{-\sin t}{2t}\to\res{-\frac12}$
  \task $t=x-\pi$: $\frac{-\sin t}{t}\to\res{-1}$
  \task $t=x-1$: $\frac{t}{t(t+2)}\to\res{\frac12}$
  \task $\tan x(\cos x-1)\sim-\frac{x^3}{2}$: \ $\res{-\frac12}$
  \task $e^{\lim\frac{-x^2/2}{x^2}}=\res{e^{-1/2}}$
\end{tasks}

\ex{2} \begin{tasks}(3)
  \task $3+\frac{\cos x}{x}\to\res{3}$
  \task $\res{0}$
  \task $t=\frac1x$: $\frac{\sin t}{t}\to\res{1}$
\end{tasks}

\ex{2} \pau{Múrcia, ordinària 2023}$\dfrac{-\frac{(2x)^2}{2}}{x\cdot x}=\res{-2}$

\ex{2} \pau{Castella i Lleó, extraordinària 2022}$\dfrac{x}{x}\to\res{1}$

\ex{2} \pau{Galícia, ordinària 2026}\begin{tasks}(2)
  \task $\frac{0}{1}=\res{0}$ (no és indeterminació)
  \task $\ln x=\ln(1+(x-1))\sim x-1$: \ $\res{1}$
  \task $\frac{1}{\cos x}\to\res{1}$
  \task $t=\frac1x\to+\infty$: $-\frac{\ln t}{\sqrt t}\to\res{0}$
\end{tasks}

\ex{2} \pau{La Rioja, extraordinària 2023}$1^\infty$: \ $e^{\lim(\cos2x-1)\frac{3}{x^2}}=e^{\lim\frac{-2x^2\cdot3}{x^2}}=\res{e^{-6}}$

\ex{3} \pau{Aragó, extraordinària 2021}$1^\infty$. Amb $t=x-1$: $\sin\left(\frac\pi2+\frac{\pi t}{2}\right)=\cos\frac{\pi t}{2}$ i $\cos\frac{\pi t}2-1\sim-\frac{\pi^2t^2}{8}$: \ $e^{\lim\frac{-\pi^2t^2/8}{t^2}}=\res{e^{-\pi^2/8}}$

\ex{2} \pau{Navarra, extraordinària 2021}$t=\frac1x\to0$: \ $\dfrac{\sin t}{t}\to\res{1}$

\ex{2} \pau{Múrcia, extraordinària 2025}$0\cdot\infty$; amb $t=\frac1x\to0^+$: \ $\dfrac{e^t-1}{t}\to\res{1}$


% =====================================================
\seccio{5}{Límits amb paràmetres}

\ex{2} \begin{tasks}(2)
  \task $\frac a2=4\Rightarrow\res{a=8}$
  \task Finit $\Rightarrow\res{a=0}$; \ $\frac b2=3\Rightarrow\res{b=6}$
  \task $e^{\lim\frac{4x(3-a)}{2x+a}}=e^{2(3-a)}\Rightarrow\res{a=2}$
  \task $\frac{mx}{2x}=\frac m2\Rightarrow\res{m=3}$
\end{tasks}

\ex{2} \begin{tasks}(1)
  \task Si $a\neq4$, $(\sqrt a-2)x\to\pm\infty$. \ $\res{a=4}$: \ $\dfrac{x}{\sqrt{4x^2+x}+2x}\to\res{\frac14}$
  \task El numerador s'ha d'anul·lar en $-3$: $9-3a+3=0\Rightarrow\res{a=4}$; \ $\dfrac{(x+3)(x+1)}{x(x+3)}\to\res{\frac23}$
\end{tasks}

\ex{2} Ha de ser $\frac00$: $1+a+b=0$; \ $x^2+ax-1-a=(x-1)(x+1+a)\Rightarrow2+a=3$: \ $\res{a=1,\ b=-2}$

\ex{2} \begin{tasks}(2)
  \task $\dfrac{(a+1)x}{2x}=\dfrac{a+1}{2}=2\Rightarrow\res{a=3}$
  \task $\dfrac{a^2x^2/2}{x^2}=\dfrac{a^2}{2}=8\Rightarrow\res{a=\pm4}$
\end{tasks}

\ex{3} \begin{tasks}(1)
  \task $a>1$: $\res{+\infty}$; \ $a<1$: $\res{-\infty}$; \ $a=1$: $\frac{x}{3x+1}\to\res{\frac13}$
  \task $1^\infty$: $e^{\lim\frac{-(ax^2+x)}{x+2}}$. \ $a>0$: $e^{-\infty}=\res{0}$; \ $a<0$: $\res{+\infty}$; \ $a=0$: $\res{e^{-1}}$
\end{tasks}

\ex{2} \pau{Andalusia, extraordinària 2022}$(\ln x)^2$ és un infinit d'ordre inferior a $x$: \ $\lim\dfrac{ax}{2x}=\dfrac a2=1\Rightarrow\res{a=2}$

\ex{3} \pau{Aragó, extraordinària 2025}$x^4+x^3-x^2-x=x(x-1)(x+1)^2$, \ i queda $x(x+1)^{2-2k}(x-1)^{1-2k}$. \ Si $k\leq0$, l'exponent $1-2k\geq1$: \ $\res{0}$. \ Si $k\geq1$, $(x-1)^{1-2k}$ té exponent negatiu i senar: $\lim_{x\to1^-}=-\infty$, $\lim_{x\to1^+}=+\infty$: \ el límit \rest{és infinit (no existeix)}

\ex{2} \pau{Castella-la Manxa, extraordinària 2022}Denominador nul en 1: $\res{b=-2}$ (el numerador $a+1\neq0$). \ $\lim_{x\to+\infty}=\frac a2=2\Rightarrow\res{a=4}$


% =====================================================
\seccio{6}{Continuïtat de funcions}

\ex{2} \begin{tasks}(1)
  \task $\lim\frac{3x}{x}=3=f(0)$: \ \rest{contínua}
  \task $\lim_{x\to0^-}=1$, $\lim_{x\to0^+}=0$: \ \rest{salt finit}
  \task $\cos\frac1x$ fitada: $\lim=0$, $\nexists h(0)$: \ \rest{evitable}; \ $\res{h(0)=0}$
\end{tasks}

\ex{2} \begin{tasks}(1)
  \task $x=0$: $f(0)=1=\lim_{x\to0^+}\frac{e^x-1}{x}$, contínua. \ $x=2$: $\lim^-=\frac{e^2-1}{2}\approx3{,}19\neq f(2)=2$: \rest{salt finit}. \ Contínua en $\res{\mathbb R-\{2\}}$
  \task $x=-2$ (és de $x<-1$): \rest{asimptòtica}. \ $x=-1$: $\lim^-=3$, $g(-1)=1$: \rest{salt finit}. \ $x=1$: $1=1$, contínua. \ Contínua en $\res{\mathbb R-\{-2,-1\}}$
  \task $x=0$: $\lim^-=2=h(0)$; \ $x=\pi$: $\lim^-=\pi+2=h(\pi)$. \ $\res{\text{Contínua en }\mathbb R}$
\end{tasks}

\ex{2} \begin{tasks}(1)
  \task $x=0$: $\lim^-=\frac{0}{1}=0$, $\lim^+=\frac{0}{+\infty}=0$, $\nexists f(0)$: \ \rest{evitable}
  \task $x=2$: $\lim^-=-(x+2)\to-4$, $\lim^+=x+2\to4$: \ \rest{salt finit}
  \task $x=1$: $\lim^-=e^{-\infty}=0$, $\lim^+=e^{+\infty}=+\infty$: \ \rest{salt infinit}
\end{tasks}

\ex{2} \pau{Castella-la Manxa, extraordinària 2024}$x=3$: $\lim^-=16$, $f(3)=\lim^+=-6$: \rest{salt finit}. \ $x=4$ (és de $x\geq3$): $\frac{8}{0^\mp}=\mp\infty$: \rest{asimptòtica}. \ Contínua en $\res{\mathbb R-\{3,4\}}$

\ex{2} \pau{Castella-la Manxa, extraordinària 2020}$x=1$: $f(1)=1$, $\lim^+=-1$: \rest{salt finit}. \ $x=2$: $\lim^-=0=\ln1=f(2)$: \rest{contínua}

\ex{2} \pau{Astúries, ordinària 2025}Dos trossos polinòmics. $x=4$: $\lim^-=2$, $f(4)=3-1=2$: \ \rest{contínua en tot el domini} $[0,+\infty)$

\ex{2} \pau{Madrid, ordinària 2023}\begin{tasks}(1)
  \task $x=-1$: $\frac13=\frac13$, contínua. $x=1$ (és de $x>-1$): $\frac{2}{0^\pm}$, \rest{asimptòtica}. Contínua en $\res{\mathbb R-\{1\}}$
  \task $1^\infty$: $\frac{x^2}{2+x^2}-1=\frac{-2}{2+x^2}$; $\lim\frac{-2(2x^2-1)}{2+x^2}=-4$: \ $\res{e^{-4}}$
\end{tasks}

\ex{2} \pau{Illes Balears, extraordinària 2024}Domini $\res{[0,1)\cup(1,+\infty)}$, on és contínua. En $x=1$: $\lim\frac{\sqrt x}{x+1}=\frac12$, $\nexists g(1)$: \ \rest{discontinuïtat evitable}

\ex{2} \pau{Castella i Lleó, extraordinària 2021}$\lim_{x\to0}\dfrac{x^2/2}{x}=0=f(0)$: \ $\res{\text{contínua en }\mathbb R}$

\ex{2} \pau{Galícia, extraordinària 2025}$\lim^-=0$, $f(0)=\frac{\ln1}{1}=0=\lim^+$: \ \rest{contínua en $x=0$}

\ex{3} \pau{Navarra, ordinària 2023}Trossos continus (en $x<1$, $1-x\neq0$; en $x\geq1$, $xe^{x+1}>0$). $f(1)=\ln e^2-2=0$. \ $\lim_{x\to1^-}\dfrac{-\sin^2(\pi x)}{1-x}$; amb $t=x-1$: $\dfrac{-\sin^2(\pi t)}{-t}\sim\dfrac{\pi^2t^2}{t}\to0$. \ $\res{\text{Contínua en }\mathbb R}$

\ex{2} \pau{Madrid, extraordinària 2022}$f(0)=0$; \ $\lim_{x\to0^+}x\ln x\overset{t=1/x}{=}\lim_{t\to+\infty}\frac{-\ln t}{t}=0$: \ \rest{contínua}


% =====================================================
\seccio{7}{Funcions a trossos amb paràmetres}

\ex{2} \begin{tasks}(2)
  \task $\lim\frac{ax}{x}=a=3$: \ $\res{a=3}$
  \task $\lim\frac{ax}{x}=a=g(0)=4$: \ $\res{a=4}$
  \task $4+2a-10=0\Rightarrow\res{a=3}$, límit $7$; \ $b+\ln1=7\Rightarrow\res{b=7}$
  \task $\frac{3x}{x}\to3$: \ $\res{a=3}$; \ $\sqrt b=3\Rightarrow\res{b=9}$
\end{tasks}

\ex{2} \begin{tasks}(1)
  \task $x=-\frac\pi2$: $2=-a+b$; \ $x=\frac\pi2$: $a+b=0$: \ $\res{a=-1,\ b=1}$
  \task $x=0$: $a=\lim\frac{3x}{x}=3$; \ $x=1$: $\ln4=2b$: \ $\res{a=3,\ b=\ln2}$
\end{tasks}

\ex{2} \begin{tasks}(2)
  \task Conjugat: $\frac14$: \ $\res{k=\frac14}$
  \task $1^\infty$: $e^{a}=e^2$: \ $\res{a=2}$
\end{tasks}

\ex{2} \pau{Andalusia, ordinària 2022}\begin{tasks}(1)
  \task $x=-1$: $-1=-a+b$; \ $x=1$: $a+b=\frac12$: \ $\res{a=\frac34,\ b=-\frac14}$
  \task Cap AV ($\frac1x$ i $\frac{x^2}{x+1}$ no fallen dins dels seus trossos). \ AH $\res{y=0}$ quan $x\to-\infty$; \ AO $\res{y=x-1}$ quan $x\to+\infty$ ($x^2=(x+1)(x-1)+1$)
\end{tasks}

\ex{2} \pau{Aragó, extraordinària 2024}$f(0)=a-1$, \ $\lim_{x\to0^+}=-b\sin\frac\pi2=-b$. \ Contínua en $\mathbb R$ si $\res{b=1-a}$; si no, salt finit en $x=0$.

\ex{2} \pau{Aragó, ordinària 2024}$\lim_{x\to0}\frac{2x}{x}=2$: contínua en $\mathbb R$ si $\res{a=2}$; si $a\neq2$, discontinuïtat evitable en $x=0$.

\ex{2} \pau{Aragó, ordinària 2022}$f(0)=e^3$; \ $\lim_{x\to0^+}(1-x)^{a/x}=e^{\lim(-x)\frac ax}=e^{-a}$: \ $\res{a=-3}$

\ex{2} \pau{Aragó, extraordinària 2021}$5-a=\frac6a\Rightarrow a^2-5a+6=0$: \ $\res{a=2\ \text{o}\ a=3}$

\ex{2} \pau{Illes Balears, ordinària 2022}$f(0)=-\frac a4$, \ $\lim_{x\to0^+}=b$: \ $\res{a=-4b}$

\ex{2} \pau{Illes Balears, extraordinària 2021}\begin{tasks}(1)
  \task \rest{Contínua} (quocient de contínues amb $2x\neq0$)
  \task $\lim\frac{ax}{2x}=\frac a2=b$: \ $\res{a=2b}$
\end{tasks}

\ex{2} \pau{Canàries, ordinària 2021}\begin{tasks}(1)
  \task El denominador s'anul·la en $x=5$: $\res{b=5}$; \ $\frac{4a-2}{7}=2$: $\res{a=4}$
  \task $f(x)=\frac{x^2+2}{x+3}$: AV $\res{x=-3}$; \ $x^2+2=(x+3)(x-3)+11$: AO $\res{y=x-3}$
\end{tasks}

\ex{2} \pau{Cantàbria, extraordinària 2026}\begin{tasks}(1)
  \task $4+a^2=4a\Rightarrow(a-2)^2=0$: \ $\res{a=2}$
  \task Cap AV. Quan $x\to-\infty$, $f\to+\infty$ i $\frac fx\to-\infty$: res. Quan $x\to+\infty$, $\frac{\ln x}{x}\to0$: AH $\res{y=4a}$
\end{tasks}

\ex{2} \pau{Cantàbria, ordinària 2025}\begin{tasks}(1)
  \task $\frac{(x-2)(x^2+2x+4)}{(x-2)(x+2)}\to\frac{12}{4}=3=e^k$: \ $\res{k=\ln3}$
  \task AV $\res{x=-2}$; \ cap AH; \ $\frac{x^2+2x+4}{x+2}=x+\frac{4}{x+2}$: AO $\res{y=x}$
\end{tasks}

\ex{2} \pau{Cantàbria, extraordinària 2020}\begin{tasks}(1)
  \task $1=\frac4\pi+a$: \ $\res{a=\dfrac{\pi-4}{\pi}}$
  \task $\lim_{x\to+\infty}=\res{a}$; \ $\lim_{x\to-\infty}\sin x$: \ $\res{\nexists}$ (oscil·la)
\end{tasks}

\ex{2} \pau{Cantàbria, extraordinària 2018}\begin{tasks}(1)
  \task $x=-2$: $0=4-2a$; \ $x=0$: $0=b$: \ $\res{a=2,\ b=0}$
  \task $f(-2)=0$, $\lim^+=4-6=-2$: \ \rest{salt finit}
\end{tasks}

\ex{2} \pau{Castella-la Manxa, extraordinària 2017}$1^\infty$: $\frac{x+1}{2x+1}-1=\frac{-x}{2x+1}$, \ $e^{\lim\frac{-1}{2x+1}}=e^{-1}$: \ $\res{k=\dfrac1e}$

\ex{2} \pau{Castella i Lleó, extraordinària 2024}Límit finit $\Rightarrow\res{a=1}$; \ $\frac{-x^2/2}{bx^2}=-\frac{1}{2b}=1$: \ $\res{b=-\frac12}$

\ex{2} \pau{Catalunya, ordinària 2026}$x=0$: $5=m^2+1\Rightarrow m=\pm2$; \ $x=2$: $(2+m)^2+1=1\Rightarrow m=-2$: \ $\res{m=-2}$

\ex{2} \pau{Extremadura, extraordinària 2023}$2+a+b=0$ i $2-a+b=5$: \ $\res{a=-\frac52,\ b=\frac12}$

\ex{2} \pau{Múrcia, extraordinària 2022}\begin{tasks}(1)
  \task $\res{0}$ (comparació d'infinits)
  \task $t=x-1$: $\frac{\ln(1+t)}{t}\to1$: \ $\res{a=1}$
\end{tasks}


% =====================================================
\seccio{8}{Teorema de Bolzano}

\ex{2} $f$ polinòmica, contínua; $f(0)=-1<0$, $f(1)=1>0$. \ $f(0{,}5)=-0{,}375$, $f(0{,}7)=0{,}043$, $f(0{,}6)=-0{,}184$: \ $c\in(0{,}6;\,0{,}7)$, \ $\res{c\approx0{,}6\ldots}$ ($c\approx0{,}682$)

\ex{2} \begin{tasks}(1)
  \task $g=x^5+x^3-1$: $g(0)=-1$, $g(1)=1$; $g(0{,}8)=-0{,}16$, $g(0{,}9)=0{,}32$: \ $\res{(0{,}8;\ 0{,}9)}$
  \task $h=e^{-x}-x$: $h(0)=1>0$, $h(1)=\frac1e-1<0$
  \task $h=\cos x-x$: $h(0)=1>0$, $h\left(\frac\pi2\right)=-\frac\pi2<0$
  \task $f(1)=0<5<6=f(2)$ i $f$ contínua: valors intermedis (o Bolzano a $x^3-x-5$)
\end{tasks}

\ex{2} \begin{tasks}(1)
  \task \rest{No}: no és contínua en $x=0$ (i de fet no s'anul·la mai)
  \task \rest{No}: no és contínua en $\frac\pi2$ (i no té cap arrel en l'interval)
  \task \rest{Sí}: $f(0)=-1<0$, $f(1)=a>0$: té una arrel en $(0,1)$
  \task \rest{No}: $f(-2)=f(2)=3>0$. Bolzano no diu res, però en té dues ($\pm1$)
\end{tasks}

\ex{2} \pau{Illes Balears, ordinària 2025}$A$ és polinòmica, contínua en $[0,500]$; $A(0)=200>0$, $A(500)=-1800<0$. Per Bolzano, \rest{existeix $c\in(0,500)$ amb $A(c)=0$}

\ex{2} \pau{Castella-la Manxa, ordinària 2023}$f$ contínua; $f(0)=-10<0$, $f(2)=28>0$: \ \rest{existeix $c\in(0,2)$ amb $f(c)=0$}

\ex{2} \pau{Castella-la Manxa, extraordinària 2017}$g(x)=\cos x-2+x$, contínua; $g(0)=-1<0$, $g(2\pi)=2\pi-1>0$: \ \rest{sí que en té}

\ex{2} \pau{Castella i Lleó, extraordinària 2022}Contínua; $f(0)=-2<0$, $f\left(\frac\pi2\right)=\frac\pi2-1>0$: \ \rest{sí, per Bolzano}

\ex{2} \pau{Castella i Lleó, extraordinària 2024}$f$ contínua, $f(0)=-1<0$, $f(1)=e-1>0$: \rest{s'anul·la en $(0,1)$}. És \rest{únic} perquè $e^x$ i $x^3$ són estrictament creixents i, per tant, també ho és la seua suma: no pot tornar a valer 0.

\ex{2} \pau{Castella i Lleó, ordinària 2017}Una corba contínua que passa d'un costat a l'altre de l'eix $OX$ el talla. \ $P(0)=-1$, $P(1)=1$: \ $\res{(0,1)}$

\ex{2} \pau{Extremadura, ordinària 2018}$f(0)=-3<0$, $f(2)=7>0$, contínua: \ \rest{hi ha una arrel en $(0,2)$}

\ex{2} \pau{Madrid, ordinària 2022}\rest{Sí}: és contínua ($x^2+1\neq0$) i $f(-1)=-\frac12<0<\frac12=f(1)$

\ex{3} \pau{Catalunya, ordinària 2019}\begin{tasks}(1)
  \task Domini $\mathbb R-\{1\}$, contínua en tot el domini; en $x=1$, $\frac10$: \ AV $\res{x=1}$
  \task En $[-2,0]$ és contínua i canvia de signe: \rest{sí}. En $[0,2]$ no és contínua (falla en 1): \rest{no podem deduir-ho}. $f(-1)=\frac72>0$: \ $\res{(-2,-1)}$
\end{tasks}

\ex{2} \pau{Galícia, ordinària 2026}\rest{És possible}: Bolzano assegura almenys un tall, no un de sol. Per exemple, $f(x)=x^3-x$ en $[-2,2]$ (talla en $-1$, $0$ i $1$).

\ex{2} \pau{Navarra, ordinària 2023}Teorema dels valors intermedis: $f$ contínua en $\left[0,\frac12\right]$, $f(0)=0<\frac34<\frac32=f\left(\frac12\right)$: \ \rest{existeix $\alpha\in\left(0,\frac12\right)$}

\ex{3} \pau{Navarra, extraordinària 2022}\begin{tasks}(1)
  \task Composició de contínues ($\ln$ està definit per a $x>0$)
  \task $f\left(\frac1e\right)=\frac{\sqrt2}{2}\approx0{,}71$, $f(e)=-\frac{\sqrt2}{2}$ i $\frac{e\sqrt2}{1-e^2}\approx-0{,}60$ està entre tots dos: \ \rest{valors intermedis}
\end{tasks}

\ex{3} \pau{Navarra, extraordinària 2021}\begin{tasks}(1)
  \task En $[6,7]$, $\sin\frac{\pi(x+1)}{4}\geq-\frac{\sqrt2}{2}$ i $2^{\frac{x-5}{2}}\geq\sqrt2$: l'argument és positiu i $f$ és composició de contínues
  \task $f(6)=\log_2\frac{\sqrt2}{2}=-\frac12<0$, $f(7)=\log_22=1>0$: \ \rest{Bolzano}
\end{tasks}


% =====================================================
\seccio{9}{Asímptotes}

\ex{2} \begin{tasks}(1)
  \task AV $\res{x=-3}$ ($+\infty$ per l'esquerra, $-\infty$ per la dreta); AH $\res{y=2}$
  \task AV $\res{x=\pm1}$ (en totes dues, $-\infty$ per l'esquerra i $+\infty$ per la dreta); $2x^3-x=(x^2-1)2x+x$: AO $\res{y=2x}$
  \task AV $\res{x=1}$ ($-\infty$ per l'esquerra, $+\infty$ per la dreta); AH $\res{y=0}$ només quan $x\to-\infty$; quan $x\to+\infty$, $\frac fx\to+\infty$: cap AO
  \task AO $\res{y=2x+\frac14}$ quan $x\to+\infty$ i $\res{y=-2x-\frac14}$ quan $x\to-\infty$; cap AV
  \task Domini $x<0$ o $x>2$. AV $\res{x=0}$ ($\lim_{x\to0^-}=-\infty$) i $\res{x=2}$ ($\lim_{x\to2^+}=+\infty$); AH $\res{y=0}$
  \task En $x=0$ el límit és 1: no hi ha AV (evitable). AO $\res{y=x}$ pels dos costats
\end{tasks}

\ex{2} \pau{Andalusia, ordinària 2021}$m=a=2$; \ $n=\lim(f(x)-2x)=\lim\frac{(b+2)x+2}{x-1}=b+2$. Passa per $(1,1)$: $2+b+2=1$: \ $\res{a=2,\ b=-3}$

\ex{2} \pau{Cantàbria, ordinària 2026}El quocient de la divisió és $x+a-a^2$: \ $a-a^2=-6\Rightarrow\res{a=3\ \text{o}\ a=-2}$

\ex{2} \pau{Cantàbria, ordinària 2026}AV $\res{x=\ln2}$ ($e^{2x}=4$ i el numerador val 1); AH $\res{y=\frac14}$ quan $x\to-\infty$ i $\res{y=0}$ quan $x\to+\infty$

\ex{2} \pau{Cantàbria, extraordinària 2018}Cal grau del numerador = grau del denominador $+1$: $\res{d=2}$; \ $x^2+1=(x-2)(x+2)+5$: \ $\res{y=x+2}$

\ex{2} \pau{Cantàbria, extraordinària 2020}$\lim\frac{x^2/2}{x}=\res{0}$: no hi ha AV. $|1-\cos x|\leq2$: $\lim_{x\to\pm\infty}f=0$: AH $\res{y=0}$

\ex{2} \pau{Canàries, extraordinària 2020}AV $\res{x=-1}$ i $\res{x=1}$; cap AH; $x^3+5x^2=(x^2-1)(x+5)+x+5$: AO $\res{y=x+5}$

\ex{2} \pau{Castella i Lleó, extraordinària 2022}Domini $\mathbb R-\{2\}$; AV $\res{x=2}$; cap AH; $x^2=(2-x)(-x-2)+4$: AO $\res{y=-x-2}$

\ex{2} \pau{Illes Balears, ordinària 2025}Domini $\res{(0,+\infty)}$. $\lim_{x\to0^+}=(+\infty)(-\infty)=-\infty$: AV $\res{x=0}$. $\lim_{x\to+\infty}=0$: AH $\res{y=0}$

\ex{2} \pau{La Rioja, extraordinària 2023}Domini $\mathbb R-\{0\}$; AV $\res{x=0}$; cap AH ni AO (el grau del numerador supera en 3 el del denominador)

\ex{2} \pau{Madrid, extraordinària 2022}$x=0$: $\lim^-=\frac{1}{0^-}=-\infty$, $f(0)=3$: \rest{salt infinit}; contínua en $\mathbb R-\{0\}$. AV $\res{x=0}$ (per l'esquerra); AH $\res{y=2}$ quan $x\to-\infty$

\ex{2} \pau{País Basc, extraordinària 2025}$x^2-3x-4=(x-4)(x+1)$: AV $\res{x=-1}$ i $\res{x=4}$; AH $\res{y=0}$

\ex{2} \pau{País Basc, ordinària 2024}$\frac{x}{(x-1)^2}$: AV $\res{x=1}$ ($+\infty$ pels dos costats); AH $\res{y=0}$; cap AO

\ex{2} \pau{Extremadura, ordinària 2018}Domini $\mathbb R-\{-1,1\}$; AV $\res{x=-1}$ i $\res{x=1}$; AH $\res{y=0}$; cap AO


\end{document}
